\section{Aula 10}

Essa aula foi mágica. Não só foi muito legal com belos teoremas,
como as provas foram fantásticas também. Infelizmente, não serei ser
capaz de fazer juz a essa aula com essas míseras notas em latex,
tudo foi com belos desenhos então tentarei meu melhor em argumentar
verbalmente o que aconteceu.

Seguimos tentando provar o teorema de Kesten, mas não pelo método
do Béla descrito na seção passada. Usamos as definições de $p^-$ e $p^+$
que críamos ser condições finitas para percolação.

Vimos a um tempo atrás, que existe um $\eta$ e um $m$ tal que se $R$ é um
retângulo $[3m]\times[m]$ e $\Pr_p(H(R)) < \eta$, então para retângulos
$R_n = [3n]\times[n]$ tínhamos $\Pr_p(H(R_n)) \to 0$ quando
$n \to \infty$. Da mesma forma, se $\Pr_p(H(R)) > 1 - \eta$,
então $\Pr_p(H(R_n)) \to 1$. Com essas coisas, havíamos definido
os pontos críticos
\[
	p^{-} := \sup\{p \in [0,1] : \exists n \quad \Pr_p(V(R_n)) < \eta\} \quad
\]
e
\[
	p^{+} := \inf\{ p \in [0,1] : \exists n\quad  \Pr_p(V(R_n)) > 1 - \eta\}.
\]

Além do mais, por Russo-Seymour-Welsh (ou o que vimos na seção passada),
$1/2 \in [p^{-}, p^{+}]$. A grande ideia é mostrar que $p^{-} =
p^{+}$ e para isso
vamos ver (assim como no Béla) que se $p^{-} < p < p^{+}$ então
\[
	\frac{d\Pr_p(V(R_n))}{dp} \gtrsim \log n
\]
Para fazer isso, obviamente usaremos Margulis-Russo [\ref{trm:russo}]
e calcularemos
\[\Ex[\# \{\text{elos pivotais para }V(R_n)\}].
\]
Curiosamente, não somaremos sobre a probabilidade de um elo ser
pivotal (como no livro),
mas veremos que, com probabilidade positiva constante, $c\log n$ elos
são pivotais. Para isso
precisaremos de stickers.

Consideramos duas percolações $\Omega_A$ e $\Omega_B$ independentes
de $\Z^2$ (funciona com qualquer modelo de
arestas independentes). Sejam $E_A$
em $\Omega_A$ e $E_B$ em $\Omega_B$ eventos e suponhamos que $E_A$ pode
ser encontrada sem necessariamente explorar todas as arestas. Vamos
considerar um novo espaço idêntico em probabilidade $\Omega_{A + B}$
onde $\Omega_A$ tem prioridade sobre $\Omega_B$. Uma aresta de $\Omega_{A+B}$
tem o valor que tem em $\Omega_A$ se foi explorada lá, se não ela
tem o valor que tem em $\Omega_B$. Definimos $E_{A+B}$ o evento onde
$E_A$ ocorre e $E_B$ acontece a menos dos valores das arestas que
determinam $E_A$.

Isso daqui requer um desenho, na aula pensamos em $E_A$ como $V(R)^*$
(estrela aqui é o dual)
e $E_B$ como achar um Mickey em $R$ no primal. Então $E_{A + B}$ é como achar um
caminho dual cruzando $R$ horizontalmente e um Mickey em $R$ (que
	pode estar cruzado pelas
arestas reveladas na descoberta do cruzamento dual). Note que com um
coupling esperto,
$\Pr_p(E_{A + B}) \geq \Pr_p(E_A)\Pr_p(E_B)$.

Vamos definir um algoritmo com stickers que com probabilidade
positiva acha muitas arestas pivotais.
Primeiramente notamos que, para uma aresta ser pivotal, deve ter um
cruzamento horizontal aberto (ao menos
possivelmente dela) usando ela e um primal vertical passando por ela
também. Nosso algoritmo
será num primeiro passo achar esse cruzamento horizontal, depois
achar o vertical e por fim
achar vários outros aneis duais cada um dando uma nova aresta pivotal.
Serei muito breve
em como fazer isso pois realmente requer desenhos.

O sketch da prova é justamente esse. No primeiro passo, exploramos em
$\Omega_A$ um caminho
mais inferior que cruza horizontalmente $[3n]\times[n/2]$, note que
isso por hipótese
acontece com probabilidade ao menos $c$ e revela arestas somente abaixo
do caminho. Se conseguimos achar esse caminho, o segundo passo é procurar
em $\Omega_B$ um caminho no dual mais a esquerda que cruza
verticalmente $[n/3]\times[n]$,
por hipótese isso acontece com probabilidade positiva $c$, revelamos
arestas somente
a esquerda do caminho e, crucialmente, considerando o sticker dos
eventos dos dois primeiros passos,
achamos com probabilidade $c^2$ uma aresta pivotal $e$. Agora consideramos
um terceiro sticker que é achar $c\log n$ cercas quadradas fechadas
centradas em $e$.
Mais precisamente, consideramos o anel $D_k = B_{3^{k+1}} - B_{3^k}$
e os eventos $C_k$ para $k < M$
que existe um caminho dual em $D_k$ que cerca o interior. Por
hipótese $\Pr_p(C_k) > c^4$.
Como temos na ordem de $M \approx \log n$ aneis contidos em $R$,
cada um independente com probabilidade $c_4$ de tá fechadinho,
com probabilidade pelo menos $1/2$ temos que $c^4M/2 $ estão fechados.
Considerando
o sticker, achamos que cada um desses aneis agora nos dá também uma
aresta pivotal. Portanto com probabilidade pelo menos $c^2/2$ achamos
na ordem de $\log n$
arestas pivotais, o que prova o resultado.

Depois vimos algo mais legal ainda. Com uma ideia de branching
process muito esperta,
o Augusto sugeriu como transformar esse $\log n$ em $n^{\eps}$
para algum $\eps$ pequeno. É muito bonito, mas um pouco complicado de
explicar sem
desenhar. Como essa amostra daqui já tá muito cabulosa escrita em
palavras, tentar guardar
mentalmente mesmo só.
